% Source: https://artofproblemsolving.com/community/c3h2125595_more_problems
% Stable identity: aops-topic-2125595-q01
Since $(30-x)^2\equiv x^2\pmod{30}$, it is enough to square $x=0,1,\dots,15$ and reduce modulo $30$:
\[0,\ 1,\ 4,\ 9,\ 16,\ 25,\ 6,\ 19,\ 4,\ 21,\ 10,\ 1,\ 24,\ 19,\ 16,\ 15 .\]
The distinct values are
\[0,\ 1,\ 4,\ 6,\ 9,\ 10,\ 15,\ 16,\ 19,\ 21,\ 24,\ 25 ,\]
twelve residues in all, with sum $\boxed{150}$.

Check with the Chinese remainder theorem: $r$ is a square modulo $30$ exactly when it is a square modulo $2$, $3$ and $5$, that is, $r\bmod3\in\{0,1\}$ and $r\bmod5\in\{0,1,4\}$ (every residue is a square modulo $2$). That gives $2\cdot2\cdot3=12$ residues, as found.

(Many number theory books count only residues relatively prime to the modulus. Under that convention the quadratic residues modulo $30$ are just $1$ and $19$, with sum $20$.)
